
\section{Funções contínuas em espaços métricos}

Nessa seção, $(M, d)$ e $(N, d')$ denotarão espaços métricos arbitrários sempre que não seja explicitamente dito o contrário.

A noção de continuidade em espaços métricos generaliza aquela aprendida para funções de uma variável real a valores reais.

Iniciaremos definindo o que é continuidade em um ponto específico.

\begin{definition}\index{continuidade!definição}
    Sejam:
    \begin{itemize}
        \item $(M, d)$ e $(N, d')$ espaços métricos,
        \item $f: M \to N$ uma função, e
        \item $p\in M$.
    \end{itemize}

    Dizemos que $f$ é \emph{contínua em $p$} se, para todo $\varepsilon > 0$, existe $\delta > 0$ tal que, para todo $x \in M$,
    \[
        \text{se } d(x, p) < \delta, \text{ então } d'(f(x), f(p)) < \varepsilon.
    \]
\end{definition}

A figura \ref{fig:continuidade} ilustra a definição de continuidade de uma função $f$ em um ponto $p$.

Notemos que isso estende diretamente a definição de continuidade pontual para funções do tipo $f: A\rightarrow \mathbb R$, em que $A\subseteq \mathbb R$.

Com isso, podemos definir o que é uma função contínua.

\begin{definition}\index{continuidade!definição}
    Seja $f: M \to N$ uma função entre espaços métricos.
    Dizemos que $f$ é \emph{contínua} se $f$ é contínua em todo ponto de $M$.
\end{definition}

Assim, a continuidade é uma propriedade local, ou seja, para verificar se uma função é contínua, basta verificar se ela é contínua em cada ponto do seu domínio.

Reforçamos, também, que o comportamento da função nas imediações de um ponto ``no ambiente'' que não esteja em seu domínio é irrelevante para verificar a continuidade.
Como acontece na definição usual de continuidade pra funções de uma variável, a função $f:\mathbb R\setminus\{0\}\rightarrow \mathbb R$ dada por $f(x) = x^{-1}$ é contínua.
O comportamento de $f$ em torno de $0$ é irrelevante, pois $0$ não pertence ao domínio de $f$.

\subimport{continuidadeRn}{figura}

Como um primeiro exemplo, tratemos da função identidade.

\begin{definition}
    Seja $A$ um conjunto qualquer.

    A função $\id_A: A\rightarrow A$ dada por $\id_A(x) = x$ para todo $x \in A$ é chamada de \emph{função identidade} de $A$.
\end{definition}

\begin{proposition}\index{função identidade}
    Seja $M$ um espaço métrico.
    A função identidade de $M$ em $M$, $\id_M$, é contínua.
\end{proposition}
\begin{proof}
    Devemos ver que para todo $p\in M$ a função $\id_M$ é contínua em $p$.
    Fixe $p\in M$ arbitrário.

    Devemos ver que para todo $\epsilon>0$ existe $\delta>0$ tal que, para todo $x\in M$,

    \[
        \text{se } d(x, p) < \delta, \text{ então } d(\id_M(x), \id_M(p)) < \varepsilon.
    \]

    Fixe $\epsilon>0$ arbitrário.
    Precisamos indicar um $\delta>0$ que satisfaça a condição acima.
    Notemos que $\id_M(x) = x$ e $\id_M(p) = p$ para todo $x\in M$.
    
    Veremos que $\delta = \epsilon$ satisfaz a condição acima.
    De fato, para todo $x \in M$, se $d(x, p) < \epsilon$, então $d(\id_M(x), \id_M(p)) = d(x, p) < \epsilon$, como desejado.

    Assim, $\id_M$ é contínua em $p$.
    Como $p\in M$ foi arbitrário, concluímos que $\id_M$ é contínua.
\end{proof}

Vejamos como isso se aplica a uma situação mais concreta.

\begin{example}
    Seja $f: \mathbb R^3\rightarrow \mathbb R^3$ dada por
    \[
        f(x, y, z) = (x, y, z).
    \]
    A função $f$ é contínua, pois $f$ é a função identidade de $\mathbb R^3$ em $\mathbb R^3$.
\end{example}


Outra classe de exemplos de funções contínuas simples são as funções constantes.

\begin{definition}\index{função constante}
    Sejam $M$ e $N$ espaços métricos.
    Uma função $f: M\rightarrow N$ é dita \emph{constante} se existe $c\in N$ tal que, para todo $x\in M$, $f(x) = c$.

    Nessas condições, dizemos que $f$ é a função constante de $M$ identicamente igual a $c$.
\end{definition}
\begin{proposition}
    Sejam:
    \begin{itemize}
        \item $M$ e $N$ espaços métricos, e
        \item $c \in N$.
    \end{itemize}

    Seja $f: M\rightarrow N$ a função de $M$ identicamente igual a $c$, ou seja, $f: M\rightarrow N$ dada por $f(x) = c$ para todo $x\in M$.

    Então $f$ é contínua.
\end{proposition}


\begin{proof}
    Devemos ver que para todo $p\in M$ a função $f$ é contínua em $p$.
    Fixe $p\in M$ arbitrário.

    Devemos ver que para todo $\epsilon>0$ existe $\delta>0$ tal que, para todo $x\in M$,
    \[
        \text{se } d(x, p) < \delta, \text{ então } d'(f(x), f(p)) < \varepsilon.
    \]

    Perceba que independente de $x$, temos $f(x) = c$ e $f(p) = c$.
    Assim, $d'(f(x), f(p)) = d'(c, c) = 0$ para todo $x\in M$.

    Por isso, devemos mostrar que para todo $\epsilon>0$ existe $\delta>0$ tal que, para todo $x\in M$,
    \[
        \text{se } d(x, p) < \delta, \text{ então } 0 < \varepsilon.
    \]

    Ora, isso é verdade para qualquer $\delta>0$.

    Fixemos $\epsilon>0$ arbitrário.
    Veremos que $\delta = 123456$ satisfaz a condição acima.
    De fato, fixe $x \in M$ tal que $d(x, p) < \delta$.
    Devemos ver que $0 < \epsilon$.
    Mas isso é verdade, pois $\epsilon$ é um número real positivo.

    Assim, $f$ é contínua em $p$.
    Como $p\in M$ foi arbitrário, concluímos que $f$ é contínua.
\end{proof}


\begin{definition}\index{projeção}
    Seja $n\geq 1$ e $i$ um inteiro com $1\leq i\leq n$.
    
    A \emph{$i$-ésima projeção} de $\mathbb R^n$ em $\mathbb R$ é a função $\pi_i: \mathbb R^n\rightarrow \mathbb R$ dada por
    \[
        \pi_i(x_1, \dots, x_n) = x_i.
    \]
\end{definition}

\begin{example}
    Para o caso $n=2$, as projeções $\pi_1$ e $\pi_2$ são dadas por $\pi_1(x, y) = x$ e $\pi_2(x, y) = y$.
\end{example}

As projeções são funções contínuas.
\begin{proposition}\index{projeção}
    Seja $n\geq 1$ e $i$ um inteiro com $1\leq i\leq n$.
    
    A $i$-ésima projeção de $\mathbb R^n$ em $\mathbb R$, ou seja, a função é contínua.
    
    Ou seja, seja $\pi_i: \mathbb R^n\rightarrow \mathbb R$ dada por $\pi_i(x_1, \dots, x_n) = x_i$.

    Então $\pi_i$ é contínua.
\end{proposition}
\begin{proof}
    Seja $p=(p_1, \dots, p_n) \in \mathbb R^n$ arbitrário e $\epsilon>0$.

    Note que, dado qualquer $x=(x_1, \dots, x_n) \in \mathbb R^n$, temos que:
    \begin{equation*}
        d(\pi_i(x), \pi_i(p)) = |x_i - p_i| \leq \sqrt{(x_1 - p_1)^2 + \cdots + (x_n - p_n)^2} = d(x, p).
    \end{equation*}

    Assim, se $d(x, p)<\epsilon$, então $d(\pi_i(x), \pi_i(p))<\epsilon$.
    Logo, tomando $\delta=\epsilon$, vemos que $\pi_i$ é contínua em $p$.

    Como $p\in \mathbb R^n$ foi arbitrário, concluímos que $\pi_i$ é contínua.
\end{proof}

Assim, por exemplo, temos que:

\begin{example}
    Seja $f: \mathbb R^2\rightarrow \mathbb R$ dada por $f(x, y) = x$.
    Então $f$ é contínua.
\end{example}

\begin{example}
    Seja $f: \mathbb R\rightarrow \mathbb R$ dada por $f(x) = x$.
    Então $f$ é contínua, pois $f$ é a projeção $\pi_1$ de $\mathbb R^1$ em $\mathbb R$.
\end{example}

\begin{example}
    Seja $f: \mathbb R^3\rightarrow \mathbb R$ dada por $f(u, v, w)=v$.
    Então $f$ é contínua, pois $f$ é a projeção $\pi_2$ de $\mathbb R^3$ em $\mathbb R$.
\end{example}

